\chapter{Introducción}
\label{anexo4:cap:introduccion}
La seguridad es un aspecto fundamental en el desarrollo del software, especialmente en aplicaciones
que se dan de cara al público. Solo en el último año, se han reportado varios incidentes de seguridad
con aplicaciones de gran de renombre, como la filtración de 49 millones de datos de usuarios de Dell
\parencites{DailyDarkWeb2024Apr}{Franceschi-Bicchierai2024May} o el robo de 560 millones de entradas
de datos de usuarios en Ticketmaster \parencite{Lawler2024Jun}.

Otro aspecto relevante es que, en la mayoría de los casos, las aplicaciones que almacenan datos de
sus usuarios deben cumplir con ciertas normativas de seguridad, transparencia y privacidad, como el
Reglamento General de Protección de Datos \parencite{EuropeanParliament2016a}. El sitio web \enquote{
GDPR Enforcement Tracker}\footnote{\url{https://www.enforcementtracker.com/}} recoge una gran cantidad
de multas impuestas a empresas por incumplimiento de la normativa. Entre algunas de ellas se encuentran
la multa de 1,200 millones de euros a Meta por parte de la Comisión Irlandesa de Protección de Datos
\parencite{DataProtectionCommission2025Jun}, la multa de 746 millones de euros a Amazon por parte de la
Comisión Nacional de Protección de Datos de Luxemburgo o la reciente multa de 3 millones de euros al
centro comercial Carrefour por parte de la Agencia Española de Protección de Datos
\parencite{AEPD:2025:PS00123:misc}.

En este contexto, el presente anexo tiene como objetivo discutir aquellos aspectos de seguridad
críticos para el correcto funcionamiento de la aplicación y se comentarán las medidas de seguridad
aplicadas para garantizar la protección de los datos en el sistema.
